\documentclass{article}
\usepackage{listings}
\usepackage{graphicx} % Required for inserting images
\usepackage{amsmath, amssymb, amsthm}
\usepackage{mathtools}
\usepackage{bm}                  % Bold math symbols

%% ---- Figures & Tables -------------------------------------
\usepackage{graphicx}
\usepackage{booktabs}            % Professional table rules
\usepackage{multirow}
\usepackage{array}
\usepackage{caption}
\usepackage{subcaption}
\usepackage{float}

\title{First Report on CL and Nectar.}
\author{Amlal El Mahrouss\\amlal@ne-app.eu}
\date{March 2026}

\begin{document}

\maketitle

\section{Introduction:}

The Context Programming paradigm as introduced in our previous work[1] has evolved since, in this paper, which serves a second report of the previous one, we will note and comment on improvements and regressions.

\section{A Reminder of Nectar's $\lambda$-Execution:}

As previously stated, we can express execution as a mathematical function:

\begin{equation}
    \Theta(x) = \lambda x.(P(x))
\end{equation}
This function must have $x$ defined in a program set $\mathbb{P}$ in which it may not be contiguous.

\section{The current implementation of Nectar:}

The current implementation of Nectar allows the following syntax:

\begin{lstlisting}
#include <GenericsLibrary/std.nhh>

const main()
{
	try {
		pthread_exit();
	}
	
	return 0;
}
\end{lstlisting}
In which this segment is completely separated from data. In order to access process data you would need a trait, such as $ThreadShared$ to retrieve the current thread information.
 
\begin{lstlisting}
#include <GenericsLibrary/std.nhh>

const main()
{
	try {
		let cur := ThreadShared->current_thread_;
	    cur->exit();
    }
	
	return 0;
}
\end{lstlisting}
This could be represented logically as:
\begin{equation}
        \operatorname{Trait}(\operatorname{C}) \coloneqq \operatorname{Rules}(\operatorname{C}) \in \mathbb{P}
\end{equation}
Such traits shall be valid in $\operatorname{C} \And$ $\Theta(\operatorname{C})$. Notice that the definition has changed from $C$ to $\mathbb{P}$. It is to express trait without a circularity condition in between.

\section{Examples of programs written in Nectar:}
\subsection{A CUDA example in Nectar:}
\begin{lstlisting}
import cudaMalloc;

const main()
{
    let ptr := 0;
    let sz := 8;
    cudaMalloc(ptr, sz);
}
\end{lstlisting}
\subsection{A Checked resource example in Nectar:}
\begin{lstlisting}
#include <GenericsLibrary/std.nhh>

const main()
{
	try {
		terminate();
	}
	
	return 0;
}
\end{lstlisting}
\subsection{A simple program in Nectar:}
\begin{lstlisting}
extern exit;

const main()
{
	let _ := exit(0);
	if (_ === 0): {
		return 0;
	}
	return 0;
}
\end{lstlisting}
\section{Motivations of Nectar:}
The paper effectively tries to address underlying issues when designing Computer Systems in data separation. That is separation data and text in a Von-Neumman architecture. Examples includes: Operating System Kernels, Program Compilers, Program Linkers.
\section{Applications of Nectar:}
Nectar applies to systems programming, specifically in unprivileged and privileged code ran Natively by an Instruction Set Architecture.
\section{Conclusion:}
The conclusion of this report, although short, is improvements in understanding how data gets separated from code, clearer formulas, and a better idea on where to head from now on.

\end{document}
